% !TeX root = ../main.tex
%
\chapter{Introduction}
\pagenumbering{arabic}
\label{sec:intro}

The past century has witnessed rapid technological progress, with computing systems among its most transformative developments.
The introduction of the internet revolutionizes how information is shared and retrieved across the world.
More recently, advancements in AI technologies significantly impact different fields, including medical, ..... .
At the foundation of these advancement lies the infrastructures that runs computation.
The core part of managing the infrastructures has always been the \gls{os} \dash\ the mediator of computation and the underlying hardware.

Every \gls{os} pursues three fundamental objectives.
First, it must expose a set of \glspl{api} that are intuitive for application developers while abstracting the complexity of underlying hardware.
Second, it must efficiently \emph{manage hardware resources}, including CPU time, memory, and storage, by allocating them to applications and reclaiming them when necessary.
Third, since the advent of multiprogramming in the 1960s, the \gls{os} has been central to ensuring system \emph{security}.
Running a malicious program can endanger the confidentiality, integrity, and availability of other applications.
Early systems such as MULTICS (MULTiplexed Information and Computing Service) tackled these challenges by introducing protection domains and privilege rings \dash\ concepts that remain foundational today.









% Interface
\hdr{The need to retrofit security into legacy abstractions}

Abstractions are the lingua franca of operating system design.
A well-crafted abstraction allows the OS to balance usability, efficiency, and security enforcement.
For instance, the \emph{process} abstraction, introduced in the 1960s~\cite{denning2022systems}, remains the cornerstone of modern operating systems.
A process represents a unit of execution that consumes CPU time and interacts with system resources.
Encapsulating related functions within processes simplifies application development (each service can be developed and debugged independently), facilitates resource management (resources can be allocated and reclaimed cleanly), and enforces isolation between mutually untrusted programs.
Another fundamental abstraction is the \emph{file}, which initially represented a block of data on disk.
In UNIX-inspired systems, this abstraction was extended so that nearly all disk I/O, networking, and interprocess communication are treated as file operations, thereby unifying system access under a single interface.
This unification also allows the OS to enforce access control policies consistently across subsystems.


Despite decades of innovation, core interfaces and abstractions in modern \glspl{os} have remained largely unchanged since the 1990s.
Over the past sixty years, numerous secure operating system architectures—such as security kernels, capability-based OSes, and microkernels—have been proposed, each promising stronger isolation and verifiable security guarantees.
However, these systems rarely achieve mainstream adoption.
Most modern applications depend on long-standing standards, particularly the \gls{posix} interface, which constrains compatibility with alternative OS architectures.
Consequently, the dominant computing infrastructures of today—such as Linux and Windows—remain monolithic systems that prioritize \gls{posix} compliance over new security paradigms.

As a result, the most practical path forward has been to \emph{retrofit} security mechanisms into existing abstractions rather than replacing them.
For instance, SELinux~\cite{selinux} extends Linux’s discretionary access control () model with mandatory access control (\gls{mac}) enforcement.
Similarly, large applications such as Chrome and Firefox retrofit finer-grained isolation mechanisms into the process abstraction to limit the damage caused by compromised third-party libraries~\cite{rlbox}.

\hdr{Protection models of legacy OSes in modern landscape}

Similarly the protection model of modern OSes has remain largely the same.
The security goals of the OS is to ensure all access to system resouces satisfy the classical properties confidentiality, integrity and availablity.

The operating system act as a \emph{reference monitor} that mediate all accesses of the system. 
Thus, it is traditionally assume to be trusted, or within the system's \gls{tcb}.

Most modern OS enforces enforces two distinct security boundaries: one between itself and the userspace processes, and the other between different processes.

The first boundary, enforced through hardware ring protection, is critical to protect the \gls{os} \gls{tcb}, ensuring that all its security enforcement are untampered.

The second boundary is the foundation for security in multiprogramming systems. 
Each process is given a dedicated adress space.
They are also given a dedicated set of resource handle (UNIX file descriptors) that they uniquely access.

Interaction  between these boundaries are strictly mediated by the OS. 
System calls allows the processes to access system services and hardware resources.
The OS validate the legitimacy of these system calls (e.g., does the process's owner have the right permission to the file it is requesting?), before granting the access to the process.

\hdr{The rise of hardware-defined trust models}

While os protection models have stagnated, hardware innovation has reshaped how systems define and enforce isolation boundaries that the legacy OS protection were never designed to  handle.
Modern systems security research now evolves along two complementary direction.
Both direction expose limitations in traditional abstractions such as processes, file descriptors, and even virtual machines, and introduce security challenges for systems designer.

\hdr{Compressed protection boundaries}

On one end, we are seeing hardware support toward \emph{compressing the protection boundaries inward}.
Motivated by how large and complex modern applications has become features such as Intel \gls{pku} enable sub-process memory isolation beyond traditional page-table mechanisms.
On ARM, Pointer Authentication (PA) and Memory Tagging (MTE) introduce pointer integrity and spatial safety semantics.
Research prototypes like \gls{cheri} extend these concepts toward full hardware-enforced memory safety and capability-based isolation~\cite{cheri}.
These mechanisms support hardware-backed partitions within a single process, enabling components to coexist under differentiated privileges without incurring the high overhead of separate processes.

Taking this principle to the extreme, CPU vendors are introducing features that enable creating secure "enclaves" that are isolated from the rest of the system, even from compromised operating systems and hypervisors.
This allow for remote cloud user to deply trusted computation with the smallest TCB. 


Hardware primitives cannot be applied superficially; incorrect integration often nullifies their security benefits.
Attacks on in-process isolation mechanisms~\cite{pku-pitfall,schrammel2022jenny,endokernel} demonstrate that existing OS interfaces must be rethought to safely support sub-process compartments.

\hdr{Expanded protection boundaries}
On the other end, we are seing the protection model \emph{expanded outward} toward whole-system protection.
Confidential Virtual Machines (CVMs), the latest development, encapsulate entire OS instances within hardware-encrypted enclaves.
This simplify the development by bring the application along with all of operational software within a single package.

CVMs, while shifting trust into hardware, inherit a large trusted computing base (TCB) from their guest OS—often millions of lines of vulnerable code.
	      Furthermore, they remain susceptible to side-channel attacks from malicious hypervisors or co-located tenants, as the OS itself may amplify information leakage through shared resources.


\hdr{Research question}
This thesis explores the following research question:

\begin{displayquote}
	\emph{How can operating systems be retrofitted to align with the trust boundaries introduced by modern hardware security mechanisms?}
\end{displayquote}

To address this question, this work explores both directions of isolation: sub-process fine-grained boundaries and system-wide trusted execution boundaries.

\begin{itemize}
	\item[-] \textbf{Retrofitting hardware-assisted fine-grained isolation into the process abstraction using capabilities.}
	      \capacity retrofits capability-based security into in-process domains~\cite{cheri,capsicum,eros,cheri-jni}.
	      Leveraging modern ARM features such as Pointer Authentication (\gls{pa}) and Memory Tagging (\gls{mte}), \capacity transforms process references into cryptographically verifiable capabilities.
	      It enforces coherent protection across the lifecycle of resources—from file import to memory access—providing hardware-backed integrity and confidentiality guarantees.
		      [[PLACEHOLDER: brief summary of \capacity’s evaluation or contribution impact]]

	\item[-] \textbf{Retrofitting OS abstractions for whole-system side-channel security.}
	      \incognitos reimagines the OS as a self-obfuscating entity within a Confidential VM.
	      It employs adaptive scheduling and memory re-randomization to obscure side-channel signals, maintaining secrecy even when the underlying hypervisor is malicious.
		      [[PLACEHOLDER: add short statement on its contribution or results]]


\end{itemize}
